\section{Introduction}
\label{sec:introduction}

European railway operators routinely deal with several thousand passenger services per day per country. Predicting which of those services will be \emph{disrupted} -- arriving late by more than five minutes or being cancelled -- before the train departs and again at every stop along its route would let operators pre-stage staff, warn passengers, and reroute when necessary. The task is a textbook example of structured spatio-temporal prediction, but it is complicated in practice by three things:

\begin{enumerate}
    \item \textbf{Data heterogeneity.} Italy, Finland and the Netherlands publish operational data in three completely different shapes: a flat per-stop CSV with sentinel-coded delays, a serialised \texttt{timeTableRows} cell carrying nested Finnish Meteorological Institute (FMI) weather observations, and a wide-tariff RDT export with three-way cancellation flags and no embedded weather. The data sources are inconsistent on the schema level, the missingness conventions, the date format, and even the meaning of the disruption label.
    \item \textbf{Class imbalance.} The fraction of stops labelled disrupted is around \SI{8}{\percent} on the held-out June test set, but skews to \SI{38}{\percent} in some countries depending on the weather window. Naive accuracy is a misleading headline metric.
    \item \textbf{Causal heterogeneity.} The Netherlands publishes a separate disruption-incident log with cause labels (rolling stock, infrastructure, accidents, weather, etc.), but Italy and Finland do not. Cross-country transfer of cause models is therefore a domain-shift problem with no in-target validation signal.
\end{enumerate}

\paragraph{Contributions.} This work contributes:
\begin{enumerate}
    \item A unified anti-leakage feature store harmonising 16.6~M stop rows across IT/FI/NL into 37 (pre-departure) or 40 (inflight) features under a single Common Data Model (\cref{sec:data,sec:features}).
    \item A four-model benchmark (Logistic Regression, LightGBM~\cite{ke2017lightgbm}, XGBoost~\cite{chen2016xgboost}, GraphSAGE~\cite{hamilton2017inductive}) trained twice -- pre-departure (A) and inflight (B) -- showing a clear inflight uplift driven by a single feature, \texttt{prev\_stop\_actual\_delay} (\cref{sec:models,sec:evaluation}).
    \item A SHAP~\cite{lundberg2017unified,lundberg2020local} TreeSHAP analysis on the full 2.30~M-row test set, plus a lag-feature ablation that shows lag history is responsible for almost the entire test-set PR-AUC of the inflight scenario (\cref{sec:xai}).
    \item An embedded Knowledge Graph -- built with Kuzu~\cite{feng2023kuzu}, designed to swap to Sparksee~\cite{sparksee2024,martinez2017dex} -- of \num{1.22}~M services, \num{2397} stations and \num{16.59}~M stop relationships that closes the loop between SHAP attributions and queryable spatial context (\cref{sec:kg}).
    \item A five-model cause-of-disruption classifier on Dutch labelled data, with a v2 redesign that lifts test macro-F1 by 46\,\% (\cref{sec:cause}). We document where transfer to Italy and Finland fails (\cref{sec:transfer}).
\end{enumerate}

The whole pipeline -- 32~unit tests, the Knowledge Graph, the model-zoo benchmark and the SHAP report -- is reproducible from raw open data with three shell commands; data ingestion and feature engineering are in pandas + scikit-learn~\cite{pedregosa2011scikit}, the deep models in PyTorch~\cite{paszke2019pytorch} and PyTorch Geometric~\cite{fey2019fast}.
